\section{Conclusion}
\label{sec:conclusion}

\subsection{Findings}

% PLACEHOLDER: replace with the verified outcome once the runs complete.
% State only what the held-out numbers support, and state simulation error
% alongside model error.
This section will report whether the discrete hazard formulation reproduces
simulated MRCs within simulation error on the held-out segment, and whether the
two parametrizations differ by more than that band. Three outcomes are
distinguishable in advance: both parametrizations track the curve within the
simulation band; the per-bin model tracks it and the shared model does not,
which would indicate that the hazard's dependence on the bin index is not
learnable from a single response surface; or neither tracks it, which would
point at the feature set rather than at the survival formulation. No claim of
the first outcome is made until \cref{tab:mrc-error} is populated.

\subsection{Limitations}

\textbf{Trace vintage.} The traces date from 2020. Cache workloads change with
application mix, object sizes, and client libraries, so the fitted hazard is a
description of that period, not a forecast. Any deployment would refit on
current traces.

\textbf{Single-vendor workload.} All clusters come from one production fleet.
Cluster-level heterogeneity is visible within the fleet, but cross-vendor
generalization is untested, and the cluster features used by the model may
encode vendor-specific conventions.

\textbf{Hazard binning.} Discretizing reuse distance into $K$ bins bounds the
resolution of the curve, and the choice of $K$ interacts with capacity: a bin
wider than the spacing of the capacity grid cannot represent the curve between
grid points. Binning also makes the estimator sensitive to the bin edges near
the knee of the curve.

\textbf{Replacement policy.} Identification relies on the LRU hit condition of
\cref{prop:lru-hit}. Production caches frequently use segmented or sampled LRU,
and their hit condition is not a pure threshold on reuse distance, so the curve
estimated here is the LRU curve for the same stream rather than the deployed
policy's curve.

\subsection{Future Work}

The natural extensions are a policy-aware recency structure to replace the
threshold condition, a continuous rather than binned hazard to remove the
resolution limit, and a per-cluster transfer check to test whether a hazard
fitted on one cluster transfers to another within the same fleet.
